\documentclass[12pt]{article}
\usepackage[utf8]{inputenc}
\usepackage[T1]{fontenc}
\usepackage{hyperref}
\title{Wahau Kenyah Language Resources for K-12 Curriculum Development}
\author{Compiled by Claude Sonnet 5.5 (claude-sonnet-5-5) for Andr\'as Kornai}
\date{}
\begin{document}
\maketitle

\section{Introduction}
Wahau Kenyah (ISO 639-3 whk; Glottocode waha1238) is an Austronesian language of East Kalimantan, Indonesia. Sources give speaker counts between about 8,000 and 10,000, and other names include Lebu' Kulit, Uma Timai and Uma' Ujok. It should not be confused with Wahau Kayan (whu), a different language of the same area. Web searches found very little published material on this language. No school primer, textbook, dictionary or grammar was found.

\section{Primary Educational Resources}
No primers, textbooks or school curricula in Wahau Kenyah turned up in search results. The language is not known to be taught in schools.

\section{Linguistic and Reference Materials}
Wikipedia has a short overview of the language: \url{https://en.wikipedia.org/wiki/Wahau_Kenyah_language}.
A conference handout from a Japanese project on language dynamics notes that serious documentation of the Kalimantan languages began only in the early 2000s, with a documentation project for Kenyah and Punan languages. It also lists Wahau Kenyah among the Kenyahic languages: \url{https://lingdy.aa-ken.jp/wp-content/uploads/2012/01/120217_soriente_inagaki_h2.pdf}.
Verbix has an endangered-languages page: \url{https://docs.verbix.com/EndangeredLanguages/WahauKenya}.
No dedicated dictionary or grammar was found.

\section{Religious and Community Resources}
Scripture Earth lists Bible access for this language: \url{https://www.scriptureearth.org/00eng.php?iso=whk}.
Joshua Project has a people-group profile for Kenyah, Wahau in Indonesia, placing the community in East Kalimantan: \url{https://joshuaproject.net/people_groups/print/12626/ID}.
No community organisation was identified.

\section{Digital and Audio Resources}
Global Recordings Network has a language page, but it says no recordings are currently available: \url{https://globalrecordings.net/en/language/whk}.

\section{Curriculum Development Context}
Very little exists for this language. A curriculum would need new documentation work with the community first: collecting texts and audio, agreeing an orthography, and building a basic word list. The related Kenyah literature could be a starting point, but it must be checked against the Wahau variety. Nothing found shows any existing school use.
\end{document}